% Source: https://artofproblemsolving.com/community/c3h2488690_vertical_asymptote_question
% Stable identity: aops-topic-2488690-q01
\textbf{(a)} For $x\ne0$, $\left|\dfrac5x\right|>M\iff|x|<\dfrac5M$. So whenever $0<|x|<\dfrac5M$, the value $\left|\dfrac5x\right|$ exceeds $M$. As $M$ can be as large as we like, $|f(x)|$ becomes arbitrarily large as $x\to0$ (from both sides: $\dfrac5x\to+\infty$ as $x\to0^+$ and $\to-\infty$ as $x\to0^-$). So $x=0$ is a vertical asymptote. $\blacksquare$

\textbf{(b)} The function is defined for $x\ne0,1$, and there
\[\dfrac{x^2-1}{x^2-x}=\dfrac{(x-1)(x+1)}{x(x-1)}=\dfrac{x+1}{x}.\]
\begin{itemize}
\item As $x\to0$, the numerator tends to $1$ and the denominator to $0$, so $|f(x)|\to\infty$: $x=0$ is a vertical asymptote.
\item As $x\to1$, $f(x)=\dfrac{x+1}{x}\to2$, so $f$ stays bounded near $1$: the graph has a hole at $(1,2)$, not an asymptote.
\end{itemize}
The only vertical asymptote is $\boxed{x=0}$.
